% % sec 4.3.4
% \subsection{Matriz de cossenos diretores}
% \label{sec:matriz_cosseno_diretor}

% Considerando dois referenciais $N$ e $B$ em rotação relativa.
% A orientação de $B$ em relação a $N$ pode ser descrita pela matriz de cossenos diretores (\gls{dcm}), denotada por $\mathbf{C}^{B/N}$.
% Seguindo a notação de \cite{wie2008space}, cada linha de $\mathbf{C}^{B/N}$ contém as projeções de um versor de $B$ nos eixos de $N$, de modo que:

% \begin{equation}
% \begin{bmatrix}
% \vec{b}_1 \\
% \vec{b}_2 \\
% \vec{b}_3
% \end{bmatrix}
% =
% \mathbf{C}^{B/N}
% \begin{bmatrix}
% \vec{n}_1 \\
% \vec{n}_2 \\
% \vec{n}_3
% \end{bmatrix}.
% \label{eq:def_dcm_ba}
% \end{equation}

% O termo $\{\vec{n}_1,\vec{n}_2,\vec{n}_3\}$ e $\{\vec{b}_1,\vec{b}_2,\vec{b}_3\}$ são bases ortonormais de $N$ e $B$, respectivamente.
% Cada elemento $C_{ij}^{B/N}$ é o cosseno do ângulo entre o $i$-ésimo versor de $B$ e o $j$-ésimo versor de $N$.

% \begin{equation}
% C_{ij}^{B/N} = \vec{b}_i \cdot \vec{n}_j.
% \end{equation}

% %%%%%%%%%%%%%%%%%%%%%%%%%%%%%%%%%%%%%%%%%%%
% A matriz $\mathbf{C}^{B/N}$ é a \gls{dcm} que realiza a mudança de base de componentes de um vetor do referencial $N$ para o referencial $B$, assumindo que $\{\vec{n}_1,\vec{n}_2,\vec{n}_3\}$ e $\{\vec{b}_1,\vec{b}_2,\vec{b}_3\}$ são bases ortonormais em $\mathbb{R}^3$.
% Por construção, suas linhas e colunas são mutuamente ortogonais e possuem norma unitária, de modo que $\mathbf{C}^{B/N}$ é uma matriz ortogonal.
% Consequentemente, $\mathbf{C}^{B/N}$ preserva produtos internos, normas e ângulos; isto é, para quaisquer vetores $\vec{x}$ e $\vec{y}$ vale $(\mathbf{C}^{B/N}\vec{x})\cdot(\mathbf{C}^{B/N}\vec{y})=\vec{x}\cdot\vec{y}$.
% Neste caso, tem-se:

% \begin{equation}
% \mathbf{C}^{B/N}\big(\mathbf{C}^{B/N}\big)^{T}
% =
% \big(\mathbf{C}^{B/N}\big)^{T}\mathbf{C}^{B/N}
% =
% \mathbf{I}_3.
% \label{eq:ortonormalidade_C}
% \end{equation}

% O termo $\mathbf{I}_3$ é a matriz identidade $3\times 3$.
% Além disso, para rotações físicas próprias, vale $\det \mathbf{C}^{B/N} = +1$.
% A relação inversa de \eqref{eq:def_dcm_ba} é escrita como:

% \begin{equation}
% \begin{bmatrix}
% \vec{n}_1 \\
% \vec{n}_2 \\
% \vec{n}_3
% \end{bmatrix}
% =
% \big(\mathbf{C}^{B/N}\big)^{-1}
% \begin{bmatrix}
% \vec{b}_1 \\
% \vec{b}_2 \\
% \vec{b}_3
% \end{bmatrix}
% =
% \big(\mathbf{C}^{B/N}\big)^{T}
% \begin{bmatrix}
% \vec{b}_1 \\
% \vec{b}_2 \\
% \vec{b}_3
% \end{bmatrix}.
% \label{eq:relacao_inversa_dcm}
% \end{equation}

% Isso se dá pois, sendo $\mathbf{C}^{B/N}$ ortogonal, sua inversa coincide com a transposta, o que corresponde à rotação inversa (transformação de $B$ para $N$).
% Esta matriz não deve ser confundida com a matriz de rotação obtida diretamente a partir de um triplo de ângulos de Euler: para um dado conjunto de ângulos de Euler existe uma DCM associada, mas a DCM fornece uma descrição independente da parametrização escolhida.
% %%%%%%%%%%%%%%%%%%%%%%%%%%%%%%%%%%%%%%%%%%%%%%%%%%%%%%%%%%%%%%%%%%%%%%%%%%%%%%%%%%%%%%%%%%%%%%%%%%%%%%%%%%
% Segundo \cite{wie2008space}, considera-se que a velocidade angular de $B$ em relação a $N$ é o vetor:

% \begin{equation}
% \vec{\omega}^{B/N}(t)
% =
% \omega_1(t)\,\vec{b}_1
% +
% \omega_2(t)\,\vec{b}_2
% +
% \omega_3(t)\,\vec{b}_3.
% \end{equation}

% Para simplificar a notação, ao longo desta subseção usa-se $\vec{\omega}\equiv\vec{\omega}^{B/N}$.
% A derivada temporal dos versores de $B$, vista a partir de $N$, é dada pela cinemática de corpo rígido:

% \begin{equation}
% \dot{\vec{b}}_i = \vec{\omega}^{B/N} \times \vec{b}_i,
% \qquad i = 1,2,3.
% \label{eq:derivada_b_i}
% \end{equation}

% Partindo de \eqref{eq:relacao_inversa_dcm}, e tomando a derivada temporal vista a partir de $N$ (isto é, $\frac{d}{dt}\big|_N$), obtém-se:

% \begin{equation}
% \frac{d}{dt}\bigg|_N
% \begin{bmatrix}
% \vec{n}_1 \\
% \vec{n}_2 \\
% \vec{n}_3
% \end{bmatrix}
% =
% \frac{d}{dt}\bigg|_N
% \left(
% \mathbf{C}^{T}
% \begin{bmatrix}
% \vec{b}_1 \\
% \vec{b}_2 \\
% \vec{b}_3
% \end{bmatrix}
% \right)
% =
% \dot{\mathbf{C}}^{T}
% \begin{bmatrix}
% \vec{b}_1 \\
% \vec{b}_2 \\
% \vec{b}_3
% \end{bmatrix}
% +
% \mathbf{C}^{T}
% \begin{bmatrix}
% \dot{\vec{b}}_1 \\
% \dot{\vec{b}}_2 \\
% \dot{\vec{b}}_3
% \end{bmatrix}.
% \label{eq:derivada_relacao_inversa}
% \end{equation}

% O termo $\frac{d}{dt}\big|_N\,\vec{n}_i=\vec{0}$, pois os versores do referencial $N$ são fixos quando a derivada é tomada em $N$.
% Para simplificar a notação, usa-se $\mathbf{C}\equiv\mathbf{C}^{B/N}$.
% Substituindo \eqref{eq:derivada_b_i} em \eqref{eq:derivada_relacao_inversa}, obtém-se:

% \begin{equation}
% \dot{\mathbf{C}}^{T}
% \begin{bmatrix}
% \vec{b}_1 \\
% \vec{b}_2 \\
% \vec{b}_3
% \end{bmatrix}
% -
% \mathbf{C}^{T}
% \begin{bmatrix}
% \vec{\omega}\times\vec{b}_1 \\
% \vec{\omega}\times\vec{b}_2 \\
% \vec{\omega}\times\vec{b}_3
% \end{bmatrix}
% =
% \begin{bmatrix}
% \vec{0} \\
% \vec{0} \\
% \vec{0}
% \end{bmatrix}.
% \end{equation}

% Para manipular o produto vetorial em forma matricial, introduz-se a matriz antissimétrica $\boldsymbol{\Omega}$ associada a $\vec{\omega}$. Esta matriz é o operador que representa o produto vetorial:

% \begin{equation}
% \vec{\omega}\times\vec{v}
% =
% \boldsymbol{\Omega}\,\vec{v},
% \qquad
% \boldsymbol{\Omega}
% =
% \begin{bmatrix}
% 0      & -\omega_3 &  \omega_2 \\
% \omega_3 & 0       & -\omega_1 \\
% -\omega_2 & \omega_1 & 0
% \end{bmatrix}.
% \label{eq:def_omega_matriz}
% \end{equation}

% De modo que $\boldsymbol{\Omega}$ é antissimétrica ($\boldsymbol{\Omega}^{T} = -\boldsymbol{\Omega}$).
% Com a definição \eqref{eq:def_omega_matriz}, a expressão anterior pode ser escrita de forma compacta como:

% \begin{equation}
% \big(\dot{\mathbf{C}}^{T} - \mathbf{C}^{T}\boldsymbol{\Omega}\big)
% \begin{bmatrix}
% \vec{b}_1 \\
% \vec{b}_2 \\
% \vec{b}_3
% \end{bmatrix}
% =
% \begin{bmatrix}
% \vec{0} \\
% \vec{0} \\
% \vec{0}
% \end{bmatrix}.
% \end{equation}

% Conduzindo-o à condição matricial:

% \begin{equation}
% \dot{\mathbf{C}}^{T} - \mathbf{C}^{T}\boldsymbol{\Omega} = \mathbf{0}.
% \end{equation}

% Tomando a transposta e usando $\boldsymbol{\Omega}^{T} = -\boldsymbol{\Omega}$, obtém-se a equação cinemática diferencial para a matriz de cossenos diretores, conforme apresentado em \cite{wie2008space}:

% \begin{equation}
% \dot{\mathbf{C}} + \boldsymbol{\Omega}\,\mathbf{C} = \mathbf{0}.
% \label{eq:cinematica_dcm}
% \end{equation}


% Escrevendo \eqref{eq:cinematica_dcm} elemento a elemento, obtêm-se as nove equações diferenciais em termos dos cossenos diretores $C_{ij}$:

% \begin{equation}
% \begin{aligned}
%     \dot{C}_{11} &= \omega_3 C_{21} - \omega_2 C_{31}, \\
%     \dot{C}_{12} &= \omega_3 C_{22} - \omega_2 C_{32}, \\
%     \dot{C}_{13} &= \omega_3 C_{23} - \omega_2 C_{33}, \\
%     \dot{C}_{21} &= \omega_1 C_{31} - \omega_3 C_{11}, \\
%     \dot{C}_{22} &= \omega_1 C_{32} - \omega_3 C_{12}, \\
%     \dot{C}_{23} &= \omega_1 C_{33} - \omega_3 C_{13}, \\
%     \dot{C}_{31} &= \omega_2 C_{11} - \omega_1 C_{21}, \\
%     \dot{C}_{32} &= \omega_2 C_{12} - \omega_1 C_{22}, \\
%     \dot{C}_{33} &= \omega_2 C_{13} - \omega_1 C_{23}.
% \end{aligned}
% \end{equation}

% Essas equações mostram explicitamente como cada cosseno diretor evolui em
% função da velocidade angular $\vec{\omega}$.
% Quando $\vec{\omega}(t)$ é conhecida como função do tempo, a orientação de $B$ em relação a $N$ pode ser obtida pela integração numérica de \eqref{eq:cinematica_dcm}.
% A condição de ortonormalidade \eqref{eq:ortonormalidade_C} é uma integral de movimento da equação cinemática: se $\mathbf{C}(0)$ é ortonormal, então a solução exata de \eqref{eq:cinematica_dcm} satisfaz $\mathbf{C}(t)\mathbf{C}^{T}(t)=\mathbf{I}$ para todo $t>0$. Na prática, essa propriedade é usada como verificação de consistência numérica e, quando necessário, como critério para re-ortonormalizar a matriz integrada.

% Por fim, a comparação entre a DCM $\mathbf{C}^{B/N}$ e a matriz de rotação obtida a partir de ângulos de Euler é fundamental para entender as diferenças entre representações de atitude. A DCM atua diretamente nas bases dos referenciais, enquanto a matriz de rotação em ângulos de Euler é construída como produto de três rotações elementares em torno de eixos prescritos.


%%%% após 20/04/2026

% sec 4.3.3
\subsection{Matriz de cossenos diretores}
\label{sec:matriz_cosseno_diretor}

A orientação da espaçonave em órbita pode ser descrita pela matriz de cossenos diretores (\gls{dcm}), a qual relaciona o referencial orbital $A$ ao referencial do corpo $B$. Nesta subseção, adota-se a matriz $\mathbf{C}^{B/A}$, cujos elementos representam os cossenos diretores entre os eixos dos referenciais $B$ e $A$ \cite{wie2008space}. 
Sendo descrito matematicamente:

\begin{equation}
\begin{bmatrix}
\vec{b}_1 \\
\vec{b}_2 \\
\vec{b}_3
\end{bmatrix}
=
\mathbf{C}^{B/A}
\begin{bmatrix}
\vec{a}_1 \\
\vec{a}_2 \\
\vec{a}_3
\end{bmatrix}.
\label{eq:def_dcm_ba_orbital}
\end{equation}

Os conjuntos $\{\vec{a}_1,\vec{a}_2,\vec{a}_3\}$ e $\{\vec{b}_1,\vec{b}_2,\vec{b}_3\}$ formam bases ortonormais associadas, respectivamente, ao referencial orbital \gls{lvlh} e ao referencial fixo ao corpo. Cada elemento da matriz é definido por:

\begin{equation}
C_{ij}^{B/A}=\vec{b}_i\cdot\vec{a}_j.
\label{eq_elemento_dcm_ba}
\end{equation}

Por construção, a matriz $\mathbf{C}^{B/A}$ é ortogonal, de modo que suas linhas e colunas são mutuamente ortogonais e possuem norma unitária, sendo:

\begin{equation}
\mathbf{C}^{B/A}\big(\mathbf{C}^{B/A}\big)^{T}
=
\big(\mathbf{C}^{B/A}\big)^{T}\mathbf{C}^{B/A}
=
\mathbf{I}_3.
\label{eq:ortonormalidade_C_ba}
\end{equation}

Em que $\mathbf{I}_3$ é a matriz identidade de ordem $3$. Além disso, para rotações próprias, vale $\det\mathbf{C}^{B/A}=+1$.
A relação inversa de \eqref{eq:def_dcm_ba_orbital} é escrita como:

\begin{equation}
\begin{bmatrix}
\vec{a}_1 \\
\vec{a}_2 \\
\vec{a}_3
\end{bmatrix}
=
\big(\mathbf{C}^{B/A}\big)^{-1}
\begin{bmatrix}
\vec{b}_1 \\
\vec{b}_2 \\
\vec{b}_3
\end{bmatrix}
=
\big(\mathbf{C}^{B/A}\big)^{T}
\begin{bmatrix}
\vec{b}_1 \\
\vec{b}_2 \\
\vec{b}_3
\end{bmatrix}.
\label{eq:relacao_inversa_dcm_ba}
\end{equation}

Essa propriedade decorre do fato de que, sendo ortogonal, a inversa da matriz coincide com a sua transposta.
No contexto orbital, a velocidade angular total da espaçonave em relação ao referencial inercial é denotada por $\vec{\omega}^{B/N}$ e escrita no referencial do corpo como:

\begin{equation}
\vec{\omega}^{B/N}(t)
=
\omega_1(t)\,\vec{b}_1
+
\omega_2(t)\,\vec{b}_2
+
\omega_3(t)\,\vec{b}_3.
\label{eq_omega_BN_corpo_dcm}
\end{equation}

Partindo da derivada temporal da relação entre os versores dos referenciais orbital e do corpo, e aplicando o teorema do transporte para referenciais em rotação relativa, obtém-se a equação diferencial matricial da \gls{dcm} \cite{wie2008space}:

\begin{equation}
\dot{\mathbf{C}}^{B/A}
+
\boldsymbol{\Omega}\,\mathbf{C}^{B/A}
-
\mathbf{C}^{B/A}\,\boldsymbol{\Omega}_{A/N}^{A}
=
\mathbf{0}.
\label{eq:cinematica_dcm_orbital}
\end{equation}

Em que $\boldsymbol{\Omega}$ é a matriz antissimétrica associada à velocidade angular total $\vec{\omega}^{B/N}$, expressa no referencial do corpo, matematicamente:

\begin{equation}
\boldsymbol{\Omega}
=
\begin{bmatrix}
0 & -\omega_3 & \omega_2 \\
\omega_3 & 0 & -\omega_1 \\
-\omega_2 & \omega_1 & 0
\end{bmatrix}.
\label{eq:def_omega_matriz_orbital}
\end{equation}

Sendo que $\boldsymbol{\Omega}_{A/N}^{A}$ é a matriz antissimétrica associada à velocidade angular do referencial orbital em relação ao referencial inercial, expressa em $A$.
Para órbita circular, a velocidade angular do referencial orbital em relação ao referencial inercial é dada pela Equação~\eqref{eq_omega_A_N}, de modo que:

\begin{equation}
\boldsymbol{\Omega}_{A/N}^{A}
=
\begin{bmatrix}
0 & 0 & n \\
0 & 0 & 0 \\
-n & 0 & 0
\end{bmatrix}.
\label{eq:omega_A_N_matriz}
\end{equation}

A Equação~\eqref{eq:cinematica_dcm_orbital} mostra que a evolução temporal da matriz de cossenos diretores depende não apenas da rotação do corpo, mas também da rotação do próprio referencial orbital. Portanto, diferentemente do caso de um corpo rígido em solo, a cinemática da atitude em órbita incorpora explicitamente o efeito da velocidade angular orbital $n$.

Escrevendo \eqref{eq:cinematica_dcm_orbital} elemento a elemento, obtêm-se nove equações diferenciais para os cossenos diretores $C_{ij}^{B/A}$.
Em particular, os termos da terceira coluna da matriz, $C_{13}$, $C_{23}$ e $C_{33}$, correspondem às componentes do versor $\vec{a}_3$ expressas no referencial do corpo e, por essa razão, aparecem explicitamente na formulação do torque de gradiente de gravidade apresentada na Subseção~\ref{sec_atitude_inercial}.

Quando a velocidade angular $\vec{\omega}^{B/N}(t)$ é conhecida como função do tempo, a orientação do corpo em relação ao referencial orbital pode ser obtida pela integração numérica de \eqref{eq:cinematica_dcm_orbital}. A condição de ortonormalidade em \eqref{eq:ortonormalidade_C_ba} permanece como critério de consistência da solução numérica e, quando necessário, pode ser utilizada para reortonormalização da matriz integrada.

Por fim, a \gls{dcm} fornece uma ligação direta entre a parametrização da atitude e os cossenos diretores que aparecem nas equações diferenciais da cinemática e na expressão do torque de gradiente de gravidade. Nas subseções seguintes, essa mesma orientação será reescrita em termos de ângulos de Euler e de quatérnions, de modo a manter consistência entre as diferentes representações adotadas neste trabalho.